%  Преамбула конспекта «Механика деформируемого твёрдого тела»
%  Компилировать pdfLaTeX (кодировка utf8, шрифты T2A).

% Latin Modern не имеет T2A-начертаний (нет t2almr.fd), поэтому кириллица идёт
% через cm-super: семейство cmr в кодировке T2A.
\usepackage[T2A]{fontenc}
\usepackage[utf8]{inputenc}
\usepackage[english,russian]{babel}

\usepackage{amsmath,amssymb,amsthm,mathtools}
\usepackage{geometry}
\geometry{a4paper,margin=2.5cm}

\usepackage{tikz}
\usetikzlibrary{arrows.meta,patterns,decorations.pathmorphing,decorations.pathreplacing,shapes.geometric,calc}
\usepackage{pgfplots}
\usepgfplotslibrary{groupplots,fillbetween}
\pgfplotsset{compat=1.18}

\usepackage[most]{tcolorbox}
\tcbuselibrary{breakable}

\usepackage{enumitem}
\usepackage{booktabs}
\usepackage{hyperref}
\hypersetup{
  colorlinks=true,
  linkcolor=blue!55!black,
  citecolor=blue!55!black,
  urlcolor=blue!55!black,
  pdftitle={Механика деформируемого твёрдого тела --- конспект}
}

\hyphenation{фор-мо-из-ме-не-ния опре-де-ля-ю-щих ре-о-ло-ги-чес-ких}

%  Окружения для теорем
\theoremstyle{plain}
\newtheorem{theorem}{Теорема}
\newtheorem{proposition}[theorem]{Утверждение}
\newtheorem{corollary}{Следствие}
\newtheorem*{corollary*}{Следствие}
\newtheorem{lemma}[theorem]{Лемма}
\newtheorem{criterion}{Признак}

\theoremstyle{definition}
\newtheorem{definition}{Определение}
\newtheorem{example}{Пример}
\newtheorem{exercise}{Упражнение}

\theoremstyle{remark}
\newtheorem*{remark}{Замечание}
\newtheorem*{reminder}{Напоминание}

\renewcommand{\qedsymbol}{$\blacksquare$}

%  Пометки конспекта
%    \begin{idea} ... \end{idea}  — пояснения, которых не было на лекции
%    \begin{warn} ... \end{warn} — неразборчиво / описка / требует сверки
\newtcolorbox{idea}[1][]{
  breakable,
  enhanced,
  colback=blue!4!white,
  colframe=blue!45!black,
  boxrule=0.4pt,
  left=6pt,right=6pt,top=5pt,bottom=5pt,
  title={Пояснение},
  fonttitle=\bfseries\small,
  coltitle=white,
  attach boxed title to top left={xshift=6pt,yshift=-2pt},
  boxed title style={colback=blue!45!black,boxrule=0pt,arc=1pt},
  #1
}

\newtcolorbox{warn}[1][]{
  breakable,
  enhanced,
  colback=orange!6!white,
  colframe=orange!65!black,
  boxrule=0.4pt,
  left=6pt,right=6pt,top=5pt,bottom=5pt,
  title={Сверить с записью},
  fonttitle=\bfseries\small,
  coltitle=white,
  attach boxed title to top left={xshift=6pt,yshift=-2pt},
  boxed title style={colback=orange!65!black,boxrule=0pt,arc=1pt},
  #1
}

% Отметка начала лекции
\newcommand{\lecture}[2]{%
  \bigskip
  \noindent\hfill\textit{\small #1\ifx\relax#2\relax\else, #2\fi}\\[-0.3em]
  \noindent\rule{\textwidth}{0.4pt}
  \medskip
}

%  Макросы
\newcommand{\R}{\mathbb{R}}
\newcommand{\N}{\mathbb{N}}
\newcommand{\eps}{\varepsilon}
\newcommand{\norm}[1]{\left\lVert #1 \right\rVert}
\newcommand{\abs}[1]{\left\lvert #1 \right\rvert}
\newcommand{\sgn}{\operatorname{sign}}
\newcommand{\dd}[2]{\frac{\partial #1}{\partial #2}}
\newcommand{\dt}[1]{\dot{#1}}
\newcommand{\tr}{\operatorname{tr}}

%  Векторы и тензоры.
%  В тетради вектор помечен одной чертой снизу, тензор 2-го ранга — двумя.
%  В конспекте и то и другое набрано полужирным: \vc — вектор, \tn — тензор.
\newcommand{\vc}[1]{\boldsymbol{#1}}
\newcommand{\tn}[1]{\boldsymbol{#1}}
\newcommand{\Eten}{\boldsymbol{E}}          % единичный тензор
\newcommand{\sig}{\boldsymbol{\sigma}}      % тензор напряжений
\newcommand{\epsT}{\boldsymbol{\varepsilon}}% тензор деформации
\newcommand{\dev}{\boldsymbol{e}}           % девиатор деформации
\newcommand{\Sdev}{\boldsymbol{S}}          % девиатор напряжений
\newcommand{\nab}{\nabla}
\newcommand{\divg}{\nab\cdot}
\newcommand{\sym}[1]{\left(#1\right)^{s}}   % симметричная часть
\newcommand{\ddotp}{\mathbin{\cdot\cdot}}   % двойное скалярное произведение

% Маркер «++ в кружке» из тетради: доказательство будет дописано позже
% по записям прошлых семестров.
\newcommand{\plusplus}{%
  \tikz[baseline=(pp.base)]{%
    \node[draw, ellipse, inner xsep=3pt, inner ysep=1pt, line width=0.4pt] (pp)
      {\footnotesize $+\,+$};}%
}

%  Реологические элементы для схем (pic-и для TikZ).
%    \pic at (x,y) {spring};    — упругий элемент (пружина, закон Гука)
%    \pic at (x,y) {damper};    — вязкий элемент (демпфер, элемент Ньютона)
%    \pic at (x,y) {stvenant};  — пластический элемент (сухое трение, Сен-Венан)
%  Каждый элемент нарисован горизонтально, с выводами в (-0.9,0) и (0.9,0).
\tikzset{
  pics/spring/.style={code={
    \draw[thick] (-0.9,0) -- (-0.55,0);
    \draw[thick, decorate, decoration={zigzag, segment length=4pt, amplitude=4pt}]
      (-0.55,0) -- (0.55,0);
    \draw[thick] (0.55,0) -- (0.9,0);
  }},
  pics/damper/.style={code={
    \draw[thick] (-0.9,0) -- (-0.4,0);
    \draw[thick] (0.3,-0.3) -- (-0.4,-0.3) -- (-0.4,0.3) -- (0.3,0.3);
    \draw[thick] (0.0,-0.26) -- (0.0,0.26);
    \draw[thick] (0.0,0) -- (0.9,0);
  }},
  pics/stvenant/.style={code={
    \draw[thick] (-0.9,0) -- (-0.3,0);
    \draw[thick] (-0.3,-0.22) rectangle (0.3,0.22);
    \draw[thick] (-0.5,0.28) -- (0.5,0.28);
    \draw[thick] (-0.5,-0.28) -- (0.5,-0.28);
    \draw[thick] (0.3,0) -- (0.9,0);
  }},
}

%  Быстрая сборка: \showfiguresfalse выключает все TikZ-рисунки.
%  Нужно, когда бесплатный Overleaf не успевает в лимит времени.
\newif\ifshowfigures
\showfigurestrue

\newcommand{\inputfig}[1]{%
  \ifshowfigures
    \input{#1}%
  \else
    \begin{center}\fbox{\small рисунок выключен: \texttt{\detokenize{#1}}}\end{center}%
  \fi
}
